AI-POWERED GEO-ADAPTIVE IMPACT-BASED PARAMETRIC INSURANCE 

        FOR SWIGGY / ZOMATO DELIVERY PARTNERS 

 

 1. PROBLEM UNDERSTANDING (1–15) 

Gig workers depend on daily earnings with no income stability. 

Income is directly proportional to active working hours. 

External disruptions reduce working efficiency or working time. 

Weather alone does not always stop work (e.g., light rain is manageable). 

Real problem is income drop, not the event itself. 

Demand fluctuations lead to idle time without earnings. 

Platform downtime completely halts earning capability. 

Traffic congestion reduces delivery throughput. 

Local restrictions (curfew/events) block movement. 

Workers often continue working even in mild disruptions. 

Existing insurance ignores short-term income loss. 

Manual claims are slow, subjective, and prone to fraud. 

No system validates whether income was actually affected. 

Workers lack real-time financial safety nets. 

Need for automatic, fair, impact-based protection. 

 

 2. CORE SOLUTION IDEA (16–30) 

Build a parametric system that insures income loss only. 

Detect disruptions using real-time data streams. 

Validate whether disruption actually impacts earnings. 

Trigger payouts only when income drop is observed. 

Eliminate manual claim filing completely. 

Use weekly pricing aligned with gig work cycles. 

Personalize policies using location + behavior. 

Integrate environmental, economic, and platform signals. 

Provide instant payouts via UPI. 

Ensure fairness using hyperlocal intelligence. 

Use WhatsApp-first UX for accessibility. 

Replace event-based triggers with impact-based triggers. 

Combine multiple signals to estimate disruption severity. 

Design system to be fraud-resistant by default. 

Build a continuously learning system. 

 

 3. PERSONA DESIGN (31–40) 

Target Swiggy / Zomato delivery partners. 

Include urban and rural riders. 

Typical income range ₹500–₹1200/day. 

Work in peak slots (lunch/dinner). 

Earnings depend on deliveries/hour. 

Riders adapt to mild disruptions (raincoat, breaks). 

Severe disruptions force them to stop working. 

Prefer zero-effort, automatic systems. 

High reliance on WhatsApp. 

Need transparency and instant trust. 

 

 4. GEO-ADAPTIVE INTELLIGENCE (41–60) 

Capture GPS and map to micro-zones (2–5 km). 

Build a location profile for each zone. 

Store historical weather data per zone. 

Track past disruption events per zone. 

Identify flood-prone areas. 

Classify zones as urban/semi-urban/rural. 

Detect infrastructure quality differences. 

Apply coastal vs inland logic. 

Cluster similar zones using AI. 

Learn zone-specific tolerance thresholds. 

Update baseline weekly. 

Detect abnormal deviations from local normal. 

Adjust rules dynamically when rider moves zones. 

Incorporate seasonal patterns into baseline. 

Maintain zone-level demand patterns. 

Learn traffic behavior patterns per zone. 

Track historical income patterns per zone. 

Generate zone risk score dynamically. 

Share learnings across similar regions. 

Enable hyperlocal fairness in payouts. 

 

 5. ENVIRONMENTAL + CONTEXTUAL MODEL (61–80) 

Monitor rainfall intensity and duration. 

Measure heat index (temperature + humidity). 

Track AQI and its trend. 

Detect extreme weather events. 

Monitor wind conditions affecting safety. 

Use forecast data for predictive alerts. 

Detect sudden environmental spikes. 

Calculate environmental severity score. 

Adjust thresholds per location. 

Combine multiple environmental factors. 

Compare current vs historical patterns. 

Detect abnormal environmental deviations. 

Incorporate time-of-day sensitivity. 

Consider rider exposure duration. 

Include road condition impact (waterlogging). 

Detect visibility conditions (heavy rain/fog). 

Adjust impact based on rider behavior. 

Combine environment + traffic impact. 

Continuously update risk model. 

Feed into disruption scoring engine. 

 

 6. IMPACT-BASED TRIGGER SYSTEM (CORE FIX) (81–100) 

Define disruption levels L0–L4. 

L0 = no impact on earnings. 

L1 = slight slowdown, minimal impact. 

L2 = noticeable drop in deliveries. 

L3 = significant income loss. 

L4 = complete halt of work. 

DO NOT trigger on event alone. 

Require validation of income impact. 

Measure order volume drop in zone. 

Measure rider activity drop. 

Measure delivery time increase. 

Combine signals into disruption score. 

Use multi-trigger fusion logic. 

Apply time-weighted disruption duration. 

Ignore short noise (<30 min events). 

Apply cooldown logic to prevent repeat payouts. 

Handle partial disruptions with partial payouts. 

Handle zone exit with prorated payouts. 

Use confidence score before triggering payout. 

Ensure payout only when earning capacity drops. 

 

 7. AI ENGINE DESIGN (101–120) 

Use ML to predict disruption risk. 

Predict next-week risk for pricing. 

Use regression for income estimation. 

Use XGBoost for premium prediction. 

Use Isolation Forest for fraud detection. 

Use time-series models for forecasting. 

Continuously retrain models weekly. 

Adapt thresholds dynamically. 

Learn rider-specific patterns. 

Learn zone-specific patterns. 

Detect anomalies in claims. 

Optimize for real-time inference. 

Ensure explainability of outputs. 

Balance accuracy vs speed. 

Use lightweight models for scalability. 

Improve predictions with more data. 

Detect behavioral deviations. 

Integrate multi-source data streams. 

Continuously refine payout accuracy. 

Maintain model fairness across regions. 

 

8. PRICING & PAYOUT MODEL (121–140) 

Weekly subscription pricing model. 

Base premium for all riders. 

Add zone risk adjustment. 

Add seasonal modifiers. 

Apply trust-based discounts. 

Cap maximum payout per day. 

Calculate payout per hour lost. 

Use severity-based multipliers. 

Use duration-based scaling. 

Allow partial payouts. 

Offer voucher-based payment option. 

Allow hybrid payment (cash + vouchers). 

Provide extra value incentives. 

Ensure affordability for riders. 

Maintain insurer profitability. 

Keep pricing transparent. 

Adjust premium weekly. 

Reflect real-time risk changes. 

Prevent adverse selection. 

Balance risk pool dynamically. 

 

 9. FRAUD DETECTION & SECURITY (141–155) 

Validate GPS during disruption. 

Detect GPS spoofing patterns. 

Cross-check with weather zone. 

Validate rider activity logs. 

Detect fake inactivity. 

Monitor claim frequency. 

Detect multiple accounts. 

Use device fingerprinting. 

Analyze behavioral anomalies. 

Use AI anomaly detection. 

Store delivery proof logs. 

Assign fraud risk score. 

Restrict high-risk users. 

Update fraud models continuously. 

Ensure system integrity. 

 

 10. GAMIFICATION & INCENTIVES (156–170) 

Assign trust score to each rider. 

Increase score with consistency. 

Reward honest usage. 

Offer premium discounts. 

Increase payout caps. 

Add leaderboard system. 

Provide reward points. 

Offer milestone bonuses. 

Introduce loyalty tiers. 

Penalize fraudulent behavior. 

Encourage safe driving. 

Reward long-term engagement. 

Promote system trust. 

Align incentives with fairness. 

Build long-term retention. 

 

11. INTERACTION & UX (171–180) 

Use WhatsApp as primary interface. 

Enable onboarding via chat. 

Send real-time alerts. 

Notify payouts instantly. 

Provide simple dashboard. 

Support multiple languages. 

Ensure low data usage. 

Keep UI minimal. 

Provide transparency in payouts. 

Enable easy understanding for riders. 

 

12. REAL-WORLD COVERAGE (181–190) 

Cover heavy rain causing slowdown. 

Cover extreme heat affecting work hours. 

Cover AQI affecting working ability. 

Cover platform downtime. 

Cover demand collapse. 

Cover traffic congestion. 

Cover restricted zones. 

Cover idle time anomalies. 

Cover overcrowded zones. 

Cover multi-factor disruptions. 

 

13. FINAL POSITIONING (191–195) 

Fully automated — zero claims. 

Geo-adaptive — location-aware. 

Impact-based — human realistic. 

Multi-factor — beyond weather. 

Intelligent — continuously learning system. 

 

FINAL ONE-LINE 

“An AI-driven, geo-adaptive income protection system that pays only when real-world disruptions actually reduce a delivery worker’s ability to earn.” 

 

 